\begin{helper}
Fix natural numbers \(p,m,t\) and a prefix tuple
\[
  pref:\{0,\ldots,m-1\}\to \mathbb F_2^p\times\mathbb F_2^p
\]
whose prefix rank is \(t\). The number of possible last pairs
\((a,b)\in\mathbb F_2^p\times\mathbb F_2^p\) such that appending the dummy pair
\((a,0)\) to \(pref\) still has rank \(t\), and such that
\[
  \langle a_{\ell+1},b\rangle=1
  \quad\text{for every zero-based }\ell<m+1
\]
in the appended first-coordinate list, is at most
\[
  2^t\,2^{p-t}.
\]
\end{helper}
\begin{proof}
Let \(U\) be the span of the first coordinates in the prefix \(pref\). Since
the prefix rank is \(t\), \(\noderef{F2BadTuplePrefixSpanCard}\) gives
\[
  |U|=2^t.
\]

Consider a counted pair \((a,b)\), and let \(ab_0\) be the tuple obtained by
appending \((a,0)\) to \(pref\). The rank of the first \(m\) entries of
\(ab_0\) is the prefix rank \(t\), by
\(\noderef{F2BadTupleRankPrefixRestriction}\), while the final rank of \(ab_0\)
is also \(t\) by the defining condition on the counted pair. Thus the rank
does not increase when \(a\) is appended. Applying
\(\noderef{F2BadTupleNonincreaseStepProductBound}\) to this last step shows
that \(a\) lies in the span of the previous first coordinates. Those previous
first coordinates are exactly the first coordinates of \(pref\), so \(a\in U\).

Now inject the counted pairs into the disjoint union over \(a\in U\) of the
possible second coordinates \(b\) satisfying the same final rank and
dot-product conditions. This map records \(a\) as an element of \(U\) and keeps
the same \(b\); it is injective because the two components recover the original
pair. For each fixed \(a\in U\), the number of possible \(b\)'s is at most
\(2^{p-t}\) by \(\noderef{F2BadTupleLastBChoicesBound}\). Therefore the total
number of counted pairs is at most
\[
  \sum_{a\in U} 2^{p-t}=|U|\,2^{p-t}=2^t\,2^{p-t}.
\]
\end{proof}
